\section{Discussion and Limitations}
\label{sec:discussion_limitations}

\subsection{Domain Collapse and Weak Discrimination on Treatment-Labelled Images}
Our empirical findings reveal fundamental failure modes in deploying closed-world vision models to post-harvest sorting. The 28.56 percentage-point degradation in top-1 accuracy (from 96.92\% to 68.36\%) demonstrates the fragility of deep convolutional backbones to unconstrained environmental shifts. In curated laboratory benchmarks, models often exploit spurious background correlations (e.g., uniform illumination, neutral mats). When deployed across external markets, mobile camera optics, ambient shadowing, and surface specularities degrade feature extraction. The severe collapse on Fresh Apples (32.12\% accuracy; 393 misclassified as Rotten Apple, 255 as Rotten Orange) coincides with visual surface variations; we present the correlation between natural surface striations, lenticel spots, and decay misclassifications strictly as an untested observational hypothesis, distinguishing qualitative visual inspection from causal evidence.

Furthermore, the breakdown of OOD detection on dataset-reported treatment-labelled produce in Cohort B (all evaluated detectors AUROC $\le 58.77\%$; 42.83\% Dual-Gate UFAR-Fresh) highlights the challenge of identifying chemical treatment status from standard RGB imagery without independent biochemical assays. On these dataset-assigned labels, the frozen convolutional backbone and post-hoc scoring functions do not achieve reliable separation from natural knowns (e.g., Mahalanobis AUROC = 47.91\%, Energy = 53.13\%). Rather than asserting general visual indistinguishability, we report strictly that the evaluated post-hoc RGB detectors fail on this repository-assigned treatment cohort. Future work could evaluate whether complementary sensing modalities improve screening under these dataset-reported condition shifts, such as: (1) near-infrared and hyperspectral spectroscopy to assess internal tissue properties, (2) volatile organic compound (VOC) headspace gas sensing to detect chemical markers \cite{sattar2024sensornet}, and (3) ambient storage and cold-chain telemetry.

\subsection{Operational Trade-offs in Decision Gating}
Our benchmark demonstrates a critical divergence between continuous ranking metrics and offline image-level acceptance and rejection outcomes. While Energy scoring achieves 81.57\% AUROC on semantic shift, prospective gating under locked validation thresholds incurs an operational Known-Class Rejection Rate (KRR) of 42.43\% across external known image records (2,270/5,350 image records rejected). If these image-level rejection rates transferred to an operational line, they could increase manual-review demand; this study did not measure commodity-level diversion, throughput, labour, or grading outcomes. While the Dual-Gate cascade suppresses false acceptance on semantic shifts, both constituent gates operate downstream of the same MobileNetV2 features and remain fundamentally bounded by visual representation capacity.

\subsection{Limitations and Threats to Validity}
To preserve research transparency, we articulate key experimental boundaries:
\begin{enumerate}
\setlength{\itemsep}{0pt}\setlength{\parsep}{0pt}\setlength{\parskip}{0pt}\setlength{\topsep}{1pt}\setlength{\partopsep}{0pt}
\item \textbf{Image Records vs.\ Specimen Independence:} Evaluated datasets lack biological fruit identifiers. Multiple records may represent rotations or burst captures of shared specimens. Although pHash sensitivity auditing across Tier 1 (excluding 56 cross-partition collisions at $r \le 8$) and Tier 2 (excluding an additional 27 intra-partition multi-label collisions) confirmed detector ranking invariance, mathematical hashing cannot prove biological independence. All bootstrap summaries are conditional on image-level resampling exchangeability.
\item \textbf{Dataset-Reported Treatment Labels without Independent Chemical Assays:} Condition labels for Cohort B and C rely on FruitVision metadata \cite{bijoy2025fruitvision}. No independent gas chromatography--mass spectrometry (GC-MS) or tissue assays were conducted by the authors. We report strictly that evaluated post-hoc vision gates exhibited weak separation on this dataset-reported cohort.
\item \textbf{Residual Within-Repository Confounding:} While source matching reduces cross-repository hardware divergence, within-repository variations (different market stalls, dates, mobile phones) remain possible.
\item \textbf{Statistical Conditioning across Cohorts:} Cohorts A--C share the FruitVision Known reference partition ($N_{\text{id}} = 4,157$). Because bootstrap resampling was executed independently per cohort, marginal confidence intervals are not paired intervals for cohort differences, and cohorts are not independent replications.
\item \textbf{Offline Evaluation without Physical Conveyance:} This study benchmarks static digital images without evaluating high-speed conveyor transport, motion blur, variable camera trigger latency, or pneumatic sorting physics.
\item \textbf{Algorithmic Implementation Variants:} As detailed in Section~\ref{sec:methodology}, the evaluated ODIN implementation optimizes the sign gradient of the maximum logit, and OpenMax fits Weibull tails to 1280-dimensional global features. Results reflect these verified implementations.
\item \textbf{Exploratory Margin Proxy Scale Dependence:} The continuous Energy--MSP linear-margin proxy (exploratory) is dominated by Energy score variance ($\sigma \approx 2.15$ vs.\ $0.18$, $\rho > 0.992$); operational decision metrics rely strictly on the scale-immune Boolean conjunction rule.
\end{enumerate}
